% =====================================================================
%  Verification Experiments  (paste into your main .tex after the
%  methodology / results sections)
%
%  Figures to upload to Overleaf  figures/  folder:
%    exp_bayesian_vs_freq/fig_rank_comparison.png
%    exp_bayesian_vs_freq/fig_cold_start_metrics.png
%    exp_synthetic/fig_attack_rank_trajectory.png
%    exp_synthetic/fig_attack_defended.png
% =====================================================================

\section{Verification Experiments}

We conduct two controlled experiments to validate key design choices
in the ranking system: Bayesian engagement smoothing and resilience to
synthetic interaction attacks.  All experiments use the same 50-image,
5-category dataset with SigLIP cosine cache, seed $= 42$, and
1\,000 simulation rounds unless stated otherwise.

% ─────────────────────────────────────────────────────────────────────
\subsection{Experiment 1: Bayesian Smoothing vs.\ Frequency Baseline}
\label{sec:exp-bayesian}

\paragraph{Motivation.}
Cold-start images (newly injected at round~501) have zero interaction
history.  The system must estimate click-through rate (CTR), like rate
(LR), and average watch time (AWT) from very few observations.  We
compare two update schemes:

\begin{itemize}
  \item \textbf{Bayesian smoothing + defense} (default): EWMA-smoothed
        counts with Bayesian priors,
        $\mathrm{CTR}_i = (\tilde V_i + c\,m) / (\tilde E_i + c)$,
        where $c{=}10$, $m{=}0.10$ act as pseudo-observations that
        regularise early estimates toward the population mean, combined
        with the velocity clamp ($k{=}2$) and trust ramp ($M{=}200$).
  \item \textbf{Frequency baseline}: sliding-window raw counts over
        the last $W{=}50$ rounds,
        $\mathrm{CTR}_i = V_i^{(W)} / E_i^{(W)}$,
        with no prior and no exponential decay.
\end{itemize}

\paragraph{Setup.}
Both runs use identical parameters—including the three-image injection
at round~501 with seeded engagement—differing \emph{only} in the
engagement update rule.

\paragraph{Results.}
Figure~\ref{fig:exp1-rank} shows the rank trajectory of the three
injected images (high, mid, low quality) under each scheme.  Under
Bayesian smoothing the high-quality image reaches rank~3 immediately
and remains stable.  Under frequency counting the same image starts at
rank~12 and requires ${\sim}100$ rounds of accumulated evidence before
converging to its deserved position.

Figure~\ref{fig:exp1-metrics} overlays CTR and LR traces: the
frequency estimates (dashed) oscillate wildly—reaching 0 or 1 in
early rounds—while the Bayesian estimates (solid) remain within a
narrow band from the first observation onward.  This confirms that
Bayesian priors are essential for cold-start stability.

\begin{figure}[htbp]
  \centering
  \includegraphics[width=\linewidth]{figures/fig_rank_comparison.png}
  \caption{Cold-start rank dynamics for injected images.
           \textbf{Left:} Bayesian smoothing with defense.
           \textbf{Right:} Frequency window ($W{=}50$).
           Bayesian priors with the trust ramp produce smooth
           convergence; frequency estimates cause prolonged rank
           instability.}
  \label{fig:exp1-rank}
\end{figure}

\begin{figure}[htbp]
  \centering
  \includegraphics[width=0.92\linewidth]{figures/fig_cold_start_metrics.png}
  \caption{CTR and LR for injected images.
           Solid lines: Bayesian; dashed: frequency.
           The frequency estimator exhibits high variance ($>0.5$
           peak-to-peak) in the first 100 rounds post-injection.}
  \label{fig:exp1-metrics}
\end{figure}

% ─────────────────────────────────────────────────────────────────────
\subsection{Experiment 2: Resilience to Synthetic Interaction Attacks}
\label{sec:exp-synthetic}

\paragraph{Motivation.}
A malicious content owner may attempt to inflate their image's ranking
by generating synthetic views and likes.  We first quantify the
baseline vulnerability, then evaluate a two-layer defense mechanism.

\paragraph{Setup.}
A single low-quality image ($q_i{=}0.01$, category ``cat'') is
injected at round~200 with average category cosine similarity.
Starting at round~220 (after 20 natural rounds), an attacker injects
$N$ fake views per round with an 80\% like ratio.  We sweep
$N \in \{5, 10, 20, 50\}$ and compare against a natural baseline
($N{=}0$).  The average organic engagement rate is approximately
$64 \times 0.3 \times 0.2 \approx 3.8$ views per image per round for
displayed images.

\paragraph{Undefended results.}
Figure~\ref{fig:exp3-rank} shows the attacked image's rank trajectory
without any defense.  Without attack, the image remains at rank~11
indefinitely.  With merely 5 fake views/round ($1.3\times$ the natural
rate), the image reaches the top-3 within 30~rounds.  At 10+ fake
views/round, top-3 is achieved within 10~rounds.  Even at the lowest
intensity, only 150~synthetic interactions suffice---a rate barely
distinguishable from an enthusiastic genuine user.

\paragraph{Defense mechanism.}
We introduce a two-layer defense:
\begin{enumerate}
  \item \textbf{Velocity clamp} ($k{=}2$): each round, any image
        whose views exceed $k \times \bar{V}$ (the population mean)
        has its engagement signals clipped to the cap.  This bounds
        the per-round impact of high-intensity attacks.
  \item \textbf{Trust ramp} ($M{=}200$): engagement signals for an
        image of age $a$ rounds are attenuated by the factor
        $\tau(a) = \min(1,\; a / M)$ before feeding into the
        EWMA/Bayesian update.  New images must accumulate engagement
        gradually; early spikes are heavily discounted.
\end{enumerate}

\paragraph{Defended results.}
Figure~\ref{fig:exp3-defended} compares undefended (left) and defended
(right) rank trajectories.  The defense increases the attack cost by
an order of magnitude: at 5~views/round the attacker now needs
140~rounds and 700~total fake interactions (vs.\ 30~rounds / 150~fake
undefended), while at 20~views/round the cost rises from 200 to
2\,600~total fake interactions.  The trust ramp is the dominant factor:
even after the velocity clamp clips high-intensity injections, the
maturity discount ensures that accumulated engagement only influences
ranking slowly as the image ages.

Despite this improvement, determined attackers with sustained effort
can still eventually reach the top-3, since any positive engagement
rate above the natural baseline will eventually dominate the Bayesian
prior.  Full mitigation would require per-user anomaly detection or
reputation-based trust scoring, which we leave to future work.

\begin{figure}[htbp]
  \centering
  \includegraphics[width=0.92\linewidth]{figures/fig_attack_rank_trajectory.png}
  \caption{Undefended attack: rank trajectory of a low-quality image
           ($q{=}0.01$) at varying synthetic interaction intensities.
           The natural baseline (grey) stays at rank~11; all attack
           intensities reach the top-3 within 10--30~rounds.}
  \label{fig:exp3-rank}
\end{figure}

\begin{figure}[htbp]
  \centering
  \includegraphics[width=\linewidth]{figures/fig_attack_defended.png}
  \caption{Undefended (left) vs.\ defended (right) attack trajectories.
           The velocity clamp ($k{=}2$) combined with trust ramp
           ($M{=}200$) delays all attacks by 100+ rounds and increases
           the required fake interactions by $3$--$13\times$.}
  \label{fig:exp3-defended}
\end{figure}
